\documentclass[10pt,a4paper]{amsart}
\usepackage[margin=19mm]{geometry}
\usepackage{amsmath,amssymb,mathtools}
\usepackage[scr=rsfs,cal=euler]{mathalfa}
\usepackage{xfrac}
\usepackage[unicode,hidelinks,bookmarks=false]{hyperref}
\newcommand{\ie}{i.e.\ }
\newcommand{\cf}{cf.\ }
\newcommand{\resp}{respectively\ }
\newcommand{\Subsection}{\subsection}
\providecommand{\nobreakdash}{\nobreak\hbox{-}}
\setlength{\emergencystretch}{2em}
\multlinegap0pt
\usepackage{bm}
\usepackage{bbm}
\usepackage{dsfont}
\usepackage{xspace}
\usepackage{cancel}
\usepackage{mathtools}
\usepackage{amsthm}
\makeatletter
\@ifundefined{thm}{\newtheorem{thm}{Thm}}{}
\@ifundefined{lem}{\newtheorem{lem}[thm]{Lemma}}{}
\makeatother
\DeclareMathOperator{\ass}{Ass}
\def\m {\mathfrak m}
\title[Associated primes of powers of edge ideals of edge-weighted ]{Associated primes of powers of edge ideals of edge-weighted trees}
\author{Jiaxin Li, Tran Nam Trung, Guangjun Zhu}
\date{Source: arXiv:2509.04774v1 (CC BY 4.0)}
\begin{document}
\maketitle

\noindent\emph{Source and licence.} Associated primes of powers of edge ideals of edge-weighted trees. Version arXiv:2509.04774v1 (CC BY 4.0), distributed under the Creative Commons Attribution 4.0 International licence, as recorded in the arXiv licence metadata for this exact version and shown on the abstract page https://arxiv.org/abs/2509.04774v1. Full rights notice: licenses/arxiv-2509.04774-CC-BY-4.0.txt.\par\smallskip

\noindent\textit{Editorial scope.} Counted unit: the Lem whose source label is \texttt{sG2} in the article (source0.tex lines 487--488, source label \texttt{sG2}), delivered together with the article's own complete original proof of it (source0.tex lines 489--576, one \texttt{proof} environment of 88 lines). No mathematics is written, repaired or re-derived: statement and proof are the article's own text line for line. Required context retained uncounted: path (lines 217--225); dec-power (lines 333--339); inc-r-assoc (lines 380--381); star (lines 388--389); sG1 (lines 400--401). Every definition, display and label of those lines is retained unchanged. Omitted from this package and not redistributed: 1--216, 226--332, 340--379, 382--387, 390--399, 402--486, 577--726. Nothing inside a retained excerpt is omitted or altered: every retained body is delivered exactly as the article's lines are written, so no omission receipt of any kind accompanies this package. Citations: the retained excerpts carry no citation command at all, so no bibliography entry of the article is transcribed and no reference is redistributed. Reference adaptation: none was needed and none was made; every reference inside the delivered unit resolves inside it, so no unresolved reference or citation is produced. Printed slips kept literal: no printed word, symbol, hypothesis or number of the counted statement or of its proof was altered. Layout and class: the article's own class is \texttt{amsart} with packages including graphicx, amsmath, amsxtra, amssymb, latexsym, amscd, amsthm, xcolor; the wrapper is the lane's pinned \texttt{amsart} editorial wrapper at 10pt on A4, which restates the theorem-like environments the retained text needs and adds the article's own macro definitions that the retained text needs (\texttt{\textbackslash ass}, \texttt{\textbackslash m}), with extra wrapper packages (bm, bbm, dsfont, xspace, cancel, mathtools). The delivered numbering is the unit's own. Assets: no retained span contains a figure environment, a graphics-inclusion command, a PDF-page inclusion, a table or a TikZ picture; no image file is copied into this package, the package declares no asset dependency and no image file is redistributed with it. Licence: version arXiv:2509.04774v1 of this article is distributed under the Creative Commons Attribution 4.0 International licence, as recorded in the arXiv licence metadata for this exact version and shown on the abstract page https://arxiv.org/abs/2509.04774v1. A full rights notice is delivered at licenses/arxiv-2509.04774-CC-BY-4.0.txt. Characters: every retained excerpt is pure ASCII, so the delivered engine, XeLaTeX with its default Latin Modern text font, reports no missing character.
\begin{lem}\label{path} Assume that $(G_\omega,v)$ is an increasing weighted tree. If $G$ is not a star  graph with a root $v$, then there is a longest path $v= v_0\to v_1\to\cdots\to v_k$ in $G$ from  $v$ such that
\begin{enumerate}
    \item $v_k$ is a leaf;
    \item if $u\in N_G(v_{k-2})$ is a non-leaf, then $\omega(v_{k-1}v_{k-2})\leqslant \omega(v_{k-2}u)$;
    \item $N_G(v_{k-1})$ has only one non-leaf $v_{k-2}$;
    \item $\omega(v_{k-1}u) \leqslant \omega(v_{k-1}v_{k-2})$ for all $u\in N_G(v_{k-1})$;
    \item $\omega(v_{k-1}v_k) \leqslant \omega(v_{k-1}u)$ for all $u\in N_G(v_{k-1})$.
\end{enumerate}
\end{lem}

\begin{lem}\label{dec-power} Let $I$ be a monomial ideal and let $x^py^q$ be a monomial in $\mathcal{G}(I)$, where $p \geqslant 1$ and $q\geqslant 1$, and $x$ and $y$ are variables. For any   $f$ in $\mathcal{G}(I)$ that satisfies 
\begin{enumerate}
    \item if  $f\ne x^py^q$, then $y\nmid f$,
    \item if  $x\mid f$, then $\deg_x(f)\geqslant p$.
\end{enumerate}
Then $(I^t \colon x^py^q)= I^{t-1}$ for all $t\geqslant 2$.
\end{lem}

\begin{lem} \label{inc-r-assoc} Let $(G_\omega,v)$ be an increasing weighted tree and let $m\geqslant\mu(v)$. Then  $\ass ((v^m,I(G_\omega))^t) \subseteq \ass ((v^m,I(G_\omega))^{t+1})$ for all $t\geqslant 1$.
\end{lem}

\begin{lem} \label{star} Let $(G_\omega,v)$ be a star graph with a root $v$ and let $m>\mu(v)$. Then  $\m\in \ass (v^m,I(G_\omega))$.
\end{lem}

\begin{lem} \label{sG1} Let $(G_\omega,v)$ be an increasing weighted tree and let $m > \mu(v)$. Then  $\m\in \ass((v^m,I(G_\omega))^t)$ for all $t\geqslant s(v,G_\omega)$+1.
\end{lem}

\begin{lem} \label{sG2} Let $(G_\omega,v)$ be an increasing weighted tree and let $m > \mu(v)$. Then $\m\in \ass((v^m,I(G_\omega))^t)$ if and only if $t\geqslant s(v,G_\omega)$+1.
\end{lem}

\begin{proof} Let $I=I(G_\omega)$. By Lemma \ref{sG1}, it suffices to show that if $\m\in\ass((v^m,I)^t)$, then $t\geqslant s(v,G_{\omega})+1$. 
 We now prove this assertion by induction on $n=|V(G)|$. 
 If  $G$ is a star graph with a root  $v$, then  $s(v,G_\omega)=0$ and the assertion follows from Lemmas \ref{inc-r-assoc} and  \ref{star}.
 
 Assume that $n > 2$ and $G$ is not a star graph. Let $k=\min\{\ell\mid \m\in \ass((v^m,I)^{\ell})\}$ and let $f$ be a monomial in $R$ such that $\m = ((v^m,I)^k \colon f)$. 
We will prove that $k\geqslant s(v,G_\omega)+1$.

By Lemma \ref{path},  there exists a longest path $v= v_0\to v_1\to\cdots\to v_{s-1}\to v_s$ in $G$ from the root $v$ such that 
\begin{enumerate}
	\item $s\geqslant 2$;
	\item $v_s$ is a leaf;
	\item if $z\in N_G(v_{s-2})$ is a non-leaf, then $\omega(v_{s-1}v_{s-2})\leqslant \omega(v_{s-2}z)$;
	\item $N_G(v_{s-1})$ has only one non-leaf $v_{s-2}$;
	\item $\omega(v_{s-1}z) \leqslant \omega(v_{s-1}v_{s-2})$ for all $z\in N_G(v_{s-1})$;
	\item $\omega(v_{s-1}v_s) \leqslant \omega(v_{s-1}z)$ for all $z\in N_G(v_{s-1})$.
\end{enumerate}
First, we will prove the following  three  claims:

{\it Claim $1$}: $(v_{s-1}v_{s})^{\omega(v_{s-1}v_{s})} \nmid f$. 

If $(v_{s-1}v_{s})^{\omega(v_{s-1}v_{s})} \mid f$, then $f = g(v_{s-1}v_{s})^{\omega(v_{s-1}v_{s})}$, where $g$ is a monomial. Together with the condition  $(6)$ and  Lemma \ref{dec-power},   this yields
\[
((v^m,I)^k \colon (v_{s-1}v_{s})^{\omega(v_{s-1}v_{s})}) = (v^m,I)^{k-1}.
\]
Therefore,
$$\m = ((v^m,I)^k \colon f) = (((v^m,I)^k\colon (v_{s-1}v_{s})^{\omega(v_{s-1}v_{s})})\colon g) = ((v^m,I)^{k-1}\colon g).$$
Hence $\m\in \ass((v^m,I)^{k-1})$. This contradicts the minimality of $k$, so  $(v_{s-1}v_{s})^{\omega(v_{s-1}v_{s})}\nmid f$, as claimed.


\medskip

{\it Claim $2$}: $\deg_{v_s}(f) = \omega(v_{s-1}v_{s})-1$ and $\deg_{v_{s-1}}(f)\geqslant \omega(v_{s-1}v_{s})$. 

Note that  $v_sf\in (v^m,I)^k$, we can write $v_{s}f$ as $v_{s}f= hf_1\cdots f_k$, where $h$ is a monomial and $f_1,\ldots,f_k\in \mathcal{G}((v^m,I))$. Since  $f\notin (v^m,I)^k$,    $v_{s}\mid f_j$ for some $j\in [k]$. Therefore,  $f_j=(v_{s-1}v_{s})^{\omega(v_{s-1}v_{s})}$, since  $v_{s}$ is a leaf of $G$.
In particular, $\deg_{v_{s-1}}(f) \geqslant \omega(v_{s-1}v_{s})$ and $\deg_{v_{s}}(f) \geqslant \omega(v_{s-1}v_{s})-1$.  By Claim $1$, $(v_{s-1}v_{s})^{\omega(v_{s-1}v_{s})}\nmid f$, which forces $\deg_{v_{s}}(f) < \omega(v_{s-1}v_{s})$, and thus $\deg_{v_{s}}(f) = \omega(v_{s-1}v_{s})-1$, as claimed.


\medskip
{\it Claim $3$}: If $s(v,G'_\omega)=s(v,G_\omega)$,  where $G'_{\omega}=G_\omega\setminus v_s$, then $k\ge s(v,G_\omega)+1$. 

Let $\m'=(z\mid z\neq v_s)$. For any $z\in \m'$,  $fz\in (v^m,I)^k$ since $\m = ((v^m,I)^k \colon f)$. 
We can write $fz$ as 
\[
fz = \gamma g_1\cdots g_k,
\]
where $\gamma$ is a monomial and $g_1,\ldots,g_k\in \mathcal{G}((v^m,I))$. Since  $z\ne v_{s}$ and by  Claim $2$,  $\deg_{v_{s}}(fz)=\omega(v_{s-1}v_{s})-1$. Therefore,  $g_i\ne (v_{s-1}v_{s})^{\omega(v_{s-1}v_{s})}$ for all $i\in [k]$. In particular, $fz\in (v^m,I')^k$, which implies that  $\m'=((v^m,I')^k\colon f)$. Therefore, $\m'\in \ass((v^m,I')^k)$. Since $|V(G')| = n-1$, by the induction hypothesis, $k\geqslant s(v,G'_{\omega})+1=s(v,G_{\omega})+1$.

\medskip
 We will prove that $k\geqslant s(v,G_{\omega})+1$ by considering the following five cases.
\begin{itemize}
     \item[(i)]  $\omega(v_{s-1}v_{s}) < \omega(v_{s-2}v_{s-1})$;
      \item[(ii)]  $\omega(v_{s-1}v_{s}) = \omega(v_{s-2}v_{s-1})$ and $L_G(v_{s-1})\setminus \{v_s\}\ne \emptyset$;
    \item[(iii)] $\omega(v_{s-1}v_{s})=\omega(v_{s-2}v_{s-1})$,  $L_G(v_{s-1})=\{v_s\}$ and $L_G(v_{s-2})= \emptyset$;
    \item[(iv)] $\omega(v_{s-1}v_{s}) = \omega(v_{s-2}v_{s-1})$,  $L_G(v_{s-1})=\{v_s\}$,  $L_G(v_{s-2})\ne \emptyset$ and $\omega(v_{s-2}u) < \omega(v_{s-2}v_{s-1})$ for some $u\in L_G(v_{s-2})$; 
    \item[(v)] $\omega(v_{s-1}v_{s}) = \omega(v_{s-2}v_{s-1})$,  $L_G(v_{s-1})=\{v_s\}$,  $L_G(v_{s-2})\ne \emptyset$ and $\omega(v_{s-2}z) \ge \omega(v_{s-2}v_{s-1})$ for all $z\in L_G(v_{s-2})$.
\end{itemize}

For the cases (i) and (ii),  we  first prove that  $s(v,G'_\omega)=s(v,G_\omega)$. Therefore, by Claim $3$,  $k\ge s(v,G_\omega)+1$.  

This is trivial if (i)  holds. If (ii)  holds, then there exists  a leaf $r\in L_G(v_{s-1})$ such that $r\ne v_s$. By the conditions (5) and (6), $\omega(v_{s-1}r)=\omega(v_{s-1}v_{s})=\omega(v_{s-2}v_{s-1})$. In particular, $s(v,G'_{\omega})=s(v,G_{\omega})$. 


For the case (iv), there exists a leaf $u\in L_G(v_{s-2})$ such that $\omega(v_{s-2}u) < \omega(v_{s-2}v_{s-1})$ and $\omega(v_{s-2}u)\leqslant \omega(v_{s-2}z)$ for all  $z\in L_G(v_{s-2})$. Together with the condition $(3)$, this yields that $\omega(v_{s-2}u)<\omega(v_{s-2}z)$ for all $z\in N_G(v_{s-2})\setminus L_G(v_{s-2})$. Therefore,  $s(v,G_\omega) = s(v, (G\setminus u)_\omega)$ and  $\omega(v_{s-2}u) \leqslant \omega(v_{s-2}z)$  for all  $z\in N_G(v_{s-2})$. 
Using the same arguments  as in   Claims $1$, $2$ and $3$, we can deduce $k\geqslant s(v,G_{\omega})+1$. 

For the cases (iii) and (v), by the condition $(3)$, we have 
\begin{equation}\label{BX1}
	\omega(v_{s-1}v_{s-2}) \leqslant \omega(v_{s-2}z) \text{ for all } z\in N_G(v_{s-2}).
	\tag{$\dag$} 
\end{equation}
Note that $s(v,G'_{\omega})=s(v,G_{\omega})-1$, and  $v_{s-1}$ is a leaf of $G'_{\omega}$  by condition $(4)$. For every $z\ne v_s$, since $\m = ((v^m,I)^k \colon f)$, $zf\in (v^m,I)^k$. Therefore,  we can write $zf$ as 
\begin{equation}\label{fz}
	fz = h g_1'\cdots g_k',
	\tag{$\ddag$} 
\end{equation}
where $h$ is a monomial and $g_1',\ldots,g_k'\in \mathcal{G}((v^m,I))$. Since  $\deg_{v_s}(zf)=\omega(v_{s-1}v_s)-1$, $g_i'\ne (v_{s-1}v_s)^{\omega(v_{s-1}v_s)}$ for all $i\in[k]$. Thus, $zf\in (v^m,I')^k$.  In particular, $\m'=((v^m,I')^k\colon f)$,
where $I'=I(G'_\omega)$. 

Substituting  $z=v_{s-1}$ into the expression  (\ref{fz}), we can obtain  that $v_{s-1}\mid g_j'$ for some $j\in[k]$, since $f\notin (v^m,I)^k$.
Note that $g_i'\ne (v_{s-1}v_s)^{\omega(v_{s-1}v_s)}$ for all $i\in [k]$, thus $g_j'=(v_{s-2}v_{s-1})^{\omega(v_{s-2}v_{s-1})}$. Therefore, $\deg_{v_{s-2}}(f)\geqslant \omega(v_{s-2}v_{s-1})$.
By Claim $2$, $v_{s-2}^{\omega(v_{s-2}v_{s-1})}v_{s-1}^{\omega(v_{s-1}v_s)}|f$. Therefore, $f$ can be written as  $f=f_1f_2$, where $f_1= v_{s-2}^{\omega(v_{s-2}v_{s-1})}\cdot v_{s-1}^{\omega(v_{s-1}v_s)}$.
Note that $v_{s-1}$ is a leaf of $G'_{\omega}$, by Lemma \ref{dec-power} and  the expression (\ref{BX1}), we have $((v^m, I')^k \colon f_1)=(v^m, I')^{k-1}$. Thus
$$
\m'= ((v^m, I')^k \colon f) = (((v^m, I')^k \colon f_1)\colon f_2) = ((v^m, I')^{k-1}\colon f_2).
$$
Therefore, $\m'\in\ass((v^m, I')^{k-1})$. By the induction hypothesis, $k-1\geqslant s(v,G'_{\omega})+1=s(v,G_{\omega})$,  implying  that $k\geqslant s(v,G_{\omega})+1$.
We complete the  proof.
\end{proof}

\end{document}
